\documentclass[11pt]{article}

\usepackage[margin=1in]{geometry}
\usepackage{amsmath}
\usepackage{graphicx}
\usepackage{tikz-cd}
\usepackage{multicol}

\setlength{\parindent}{0pt}
\setlength{\parskip}{\baselineskip}

\title{Sustainable Life Management}
\author{}
\date{} 

\begin{document}
\maketitle

\section*{Problem Description}

We study a task allocation problem motivated by sustainable life management and knowledge work: within a finite planning horizon, a team must decide which tasks to undertake and how to assign them across multiple agents. Each agent has limited resources, including not only available time but also cognitive capacity and attention. Tasks are and uncertain, for example, even when a task is estimated to take one hour, its realized requirement may be substantially longer due to interruptions, switching costs, or unforeseen complexity. Meanwhile, the utility generated by completing a task is also uncertain and can depend on the realized workload and cognitive strain.

Formally, we model the system as a \emph{multiple knapsack} problem in which each agent corresponds to a knapsack with capacity $C_j$, and each task corresponds to an item. For any task--agent pair $(i,j)$, the task consumes a \emph{stochastic size} $S_{ij}$ (e.g., realized time/effort/attention usage) when processed by agent $j$, and yields a \emph{stochastic return} $R_{ij}$ (utility, productivity gain, or value). $R_{ij}$ may be correlated with $S_{ij}$ through stress--performance mechanisms: deeper engagement can increase value for some tasks, whereas overload may reduce performance and diminish returns for others.

The planner chooses binary assignment decisions $x_{ij}\in\{0,1\}$, indicating whether task $i$ is assigned to agent $j$. Each task can be assigned to at most one agent. Because resource consumption is uncertain and cognitive overload is costly, feasibility is enforced through a risk control requirement: for each agent, the probability that total realized workload exceeds capacity is restricted to be small, which is a chance constraint. The objective is to maximize total expected utility across agents subject to assignment and overload-risk constraints.

This formulation captures three features of cognitive work allocation: (i) agents with different capacities and task efficiencies, (ii) uncertain time/attention requirements that drive schedule overruns and stress, and (iii) uncertain and workload-dependent returns that reflect variable productivity under cognitive constraints. This problem generalizes classical deterministic multiple knapsack models by introducing stochastic sizes and stochastic returns with operationalizing cognitive/attention limits via constraints.


\end{document}
